% Write your conclusions here
\label{Conclusions} % If you need a label for this chapter

% Example text



This dissertation investigated low-bit post-training quantization of
Llama-3.2-3B-Instruct using AWQ, SmoothQuant, a project-specific W4A8
pipeline and NestQuant. The study focused on the trade-off between model
quality, memory reduction and practical deployment.

The results showed that no single method performed best across all
objectives. SmoothQuant achieved the strongest quality retention, with
only a 0.48\% perplexity increase relative to BF16, while NestQuant
W4.25A16KV16 followed closely at 0.64\%. AWQ and W4A8 produced larger
quality degradation but achieved the strongest memory savings, reducing
runtime model storage from 5.98~GiB to 2.13~GiB and peak GPU memory from
6.33~GiB to 2.47~GiB.

The NestQuant ablation showed that activation quantization was more
quality-sensitive than KV-cache quantization. A4 increased perplexity by
2.941\%, compared with only 0.743\% for KV4. The packed NestQuant
W4.25A16KV16 implementation also reduced checkpoint size from 5.78~GiB
to 2.16~GiB and peak GPU memory from 6.01~GiB to 2.84~GiB.

Overall, the results demonstrate that low-bit quantization should be
evaluated using model quality, quantization scope, memory reduction and
deployment practicality rather than bit width alone. SmoothQuant provided
the best quality retention, packed W4 configurations delivered the
largest memory savings, and NestQuant provided the broadest quantization
scope across weights, activations and the KV cache.
